\documentclass[11pt,aspectratio=169]{beamer}
% WiL progress update, 7 November 2026.  Same conventions as update_221026:
% one claim per slide, the slide title IS the claim, and the data that shows it
% sits on the same slide.  Every figure is the project's own output.
\usetheme{default}
\usecolortheme{seagull}
\setbeamertemplate{navigation symbols}{}
\setbeamertemplate{footline}[frame number]
\setbeamerfont{frametitle}{size=\large,series=\bfseries}
\setbeamercolor{frametitle}{fg=black!85}
\usepackage[T1]{fontenc}
\usepackage[utf8]{inputenc}
\usepackage{booktabs}
\usepackage{siunitx}
\usepackage{graphicx}
\usepackage{xcolor}
\sisetup{detect-all,range-phrase=--,range-units=single}
\newcommand{\repo}[1]{\texttt{\detokenize{#1}}}
\newcommand{\FIG}{../../final_report/figures}
\newcommand{\OUT}{../../../../output/compare/simulation}
\newcommand{\OUTR}{../../../../output/compare/real}
\newcommand{\src}[1]{\vfill{\tiny\color{black!55}Source: #1}}
\definecolor{neg}{HTML}{B00020}
\definecolor{pos}{HTML}{1B5E20}
\newcommand{\bad}[1]{\textcolor{neg}{#1}}
\newcommand{\good}[1]{\textcolor{pos}{#1}}

\title{WiL update, 7 November 2026}
\subtitle{Dataset decision \textperiodcentered\ The loop closer \textperiodcentered\
          A real non-visual baseline \textperiodcentered\ Detector \textperiodcentered\ Maps}
\author{Bolton Athit Davies}
\date{}

\begin{document}

\begin{frame}
  \titlepage
\end{frame}

\begin{frame}{Contents}
  \begin{enumerate}
    \item Dataset decision: what the report is now built from
    \item Experiment 3: the loop closer decides ORB-SLAM3's error
    \item Experiment 3: what each pipeline costs
    \item Experiment 1: a non-visual baseline on the real robot
    \item Experiment 2: the detector on the realistic worlds
    \item Experiment 4: four maps, all failing the criterion
    \item Where the report stands, and what is open
  \end{enumerate}
\end{frame}

% ===================================================================== Part 1
\section{Dataset decision}

\begin{frame}{1. The dataset set is fixed: eleven recordings, everything else retired}
  \begin{columns}[T]
    \column{0.56\textwidth}
      \scriptsize\setlength{\tabcolsep}{3pt}
      \begin{tabular}{lll}
        \toprule
        Recording & Role & Layout \\
        \midrule
        \repo{realistic_000} & Exp.\ 1, 3; Exp.\ 4 probe & U static \\
        \repo{realistic_001} & Exp.\ 1, 3; Exp.\ 4 probe & I static \\
        \repo{realistic_002} & Exp.\ 1, 3 & U traffic \\
        \repo{realistic_003} & Exp.\ 1, 3 & I traffic \\
        \repo{realistic_map} & Exp.\ 4 build; Exp.\ 3 cost & U static \\
        \repo{realistic_01_map} & Exp.\ 4 build; Exp.\ 3 cost & I static \\
        \repo{map_tiny} & Exp.\ 4 build & tiny \\
        \repo{realv6_000} & Exp.\ 1, 3 (real) & Gensurv \\
        \repo{realv7_000}--\repo{002} & Exp.\ 1, 3 (real) & Gensurv \\
        \bottomrule
      \end{tabular}
    \column{0.44\textwidth}
      \footnotesize
      \begin{itemize}
        \item Retired on 6 November: the \repo{allsensor} and
              \repo{nofloortexture} campaign, \repo{relogged_20260916}, the
              October filtered runs, \repo{realv2}--\repo{realv5}
        \item Their results went to \repo{trash/output/}, logged, and are
              cited nowhere
        \item \textbf{Dynamic environment} is now a \emph{condition} of a
              recording, not a factor: no two recordings share a route
        \item Filtered arms are blank in both domains until the repeat is run
      \end{itemize}
  \end{columns}
  \src{\repo{trash/output/RETIRED_20261106.txt}; registry in chapter~3}
\end{frame}

\begin{frame}{1. Six of eight result slots now hold data}
  \centering\small\setlength{\tabcolsep}{4pt}
  \begin{tabular}{p{0.26\textwidth}p{0.33\textwidth}p{0.33\textwidth}}
    \toprule
    Experiment & Simulation & Real robot \\
    \midrule
    1. Non-visual baseline & \good{completed}: 4 recordings, EKF at the VSLAM
      noise model & \good{partial}: wheel+gyro, fitted scale \\
    2. YOLO detector & \good{completed}: 4 recordings, 4-class checkpoint &
      planned: no annotated frames \\
    3. VIO + filtering & \good{partial}: 20 unfiltered runs + 18 cost repeats;
      \bad{no filtered run} & \good{partial}: ORB-SLAM3 provisional;
      \bad{VINS no valid run} \\
    4. Dense map & \good{partial}: tiny passes, 4 realistic
      \bad{fail}; large \bad{failed} & planned \\
    \bottomrule
  \end{tabular}
  \par\vspace{0.8em}
  \footnotesize Changed since 22 October: Experiment~2 moved from untested to
  completed, Experiment~1 gained its real-robot half, and Experiment~3 gained
  both its cost measurement and a mechanism for ORB-SLAM3's error.
  \src{\repo{chapter_4_experiments.tex}, table 4.1}
\end{frame}

% ===================================================================== Part 2
\section{Experiment 3: the loop closer}

\begin{frame}{2. Twenty unfiltered runs: VINS is steady, ORB-SLAM3 is bimodal}
  \begin{columns}[T]
    \column{0.40\textwidth}
      \footnotesize
      \begin{itemize}
        \item Factors: estimator, front end (CPU / GPU), loop closer
              (on / off). One run per cell, replay rate 1.0
        \item VINS-Fusion \SIrange{0.22}{0.66}{\metre}, no jump, on every
              recording
        \item ORB-SLAM3 \SIrange{0.14}{6.02}{\metre}: eleven runs of sixteen
              between \SI{0.14}{} and \SI{0.52}{\metre}, five diverge
        \item Neither front end ever loses tracking: zero loss events, zero
              resets in twenty runs
      \end{itemize}
    \column{0.60\textwidth}
      \includegraphics[width=\textwidth,height=0.80\textheight,keepaspectratio]{\FIG/realistic_metrics.png}
  \end{columns}
  \src{\repo{output/compare/simulation/exp-realistic-ab/runs_metrics.csv}}
\end{frame}

\begin{frame}{2. Every ORB-SLAM3 divergence is the loop closer, in both directions}
  \centering\scriptsize\setlength{\tabcolsep}{4pt}
  \begin{tabular}{llrr>{\raggedright\arraybackslash}p{0.34\textwidth}}
    \toprule
    Recording & Layout & Closer on & Closer off & What the closer did \\
    \midrule
    \repo{realistic_000} & U static & 0.34 / 0.29 & 0.48 / 0.52 &
      two closures of \SIrange{0.07}{0.12}{\metre}: inert \\
    \repo{realistic_001} & I static & \bad{6.02 / 4.78} & 0.50 / 0.40 &
      \bad{one false closure of \SI{14}{\metre}} \\
    \repo{realistic_002} & U traffic & 1.39 / 0.40 & 0.29 / 0.17 &
      one closure of \SI{0.26}{\metre}; the 1.39 is an init offset \\
    \repo{realistic_003} & I traffic & \good{0.14 / 0.50} & 1.76 / 1.91 &
      \good{one closure removes \SI{5}{\metre} of vertical drift} \\
    \bottomrule
  \end{tabular}
  \par\vspace{0.6em}
  \scriptsize Each cell is aligned ATE in metres, CPU front end / GPU front
  end. The closer state is read from each run's own log, never from its name.
  \src{\repo{loop_closing.csv} per run; \repo{exp-realistic-ab}}
\end{frame}

\begin{frame}{2. The I layout's repeated aisles produce a false 14 m closure}
  \begin{columns}[T]
    \column{0.36\textwidth}
      \footnotesize
      \begin{itemize}
        \item Both front ends accept \textbf{one} closure at
              \SI{73.2}{\second}, \SI{76}{\percent} along the route
        \item Correction \SI{14.7}{\metre} (CPU) and \SI{13.3}{\metre} (GPU)
        \item The estimate had drifted about \SI{0.5}{\metre} before it: no
              real revisit needs a \SI{14}{\metre} shift
        \item Matched keyframes are \textbf{one aisle apart}
        \item Four later candidates were rejected by the inertial geometry
              check, after the damage
        \item Both front ends failing at the same frame points at place
              recognition, not the extractor
      \end{itemize}
    \column{0.64\textwidth}
      \includegraphics[width=\textwidth,height=0.84\textheight,keepaspectratio]{\FIG/timeline_realistic__realistic_001.png}
  \end{columns}
  \src{\repo{timeline_realistic__realistic_001.png}; step at \SI{73}{\second}}
\end{frame}

\begin{frame}{2. On the same layout with traffic, only a closure removes the vertical drift}
  \begin{columns}[T]
    \column{0.36\textwidth}
      \footnotesize
      \begin{itemize}
        \item Closure-free runs read \SI{1.76}{} and
              \SI{1.91}{\metre} ATE, yet their \emph{planar} error never
              exceeds \SI{0.7}{\metre}
        \item The difference is vertical: they climb to \SI{4.7}{} and
              \SI{5.1}{\metre} above the floor
        \item The loop-closing runs stay within \SI{1.4}{\metre} and read
              \SI{0.14}{} and \SI{0.50}{\metre}
        \item Our ATE is three-dimensional, so it charges that climb to a
              floor robot: the planar onset table is the complement
      \end{itemize}
    \column{0.64\textwidth}
      \includegraphics[width=\textwidth,height=0.84\textheight,keepaspectratio]{\FIG/timeline_realistic__realistic_003.png}
  \end{columns}
  \src{No run crosses its planar threshold on this recording}
\end{frame}

\begin{frame}{2. The GPU front end changes how many frames survive, not accuracy}
  \begin{columns}[T]
    \column{0.52\textwidth}
      \footnotesize
      \begin{itemize}
        \item GPU extractor keeps \SIrange{97}{99}{\percent} of each bag's
              frames; CPU keeps \SIrange{72}{94}{\percent}
        \item Tracking median \SIrange{30}{37}{\milli\second} against
              \SIrange{49}{55}{\milli\second}
        \item Across the eight CPU/GPU pairs at the same closer state, five
              differ by at most \SI{0.16}{\metre}
        \item The three larger differences are each one run's discrete event,
              and go in both directions
        \item \textbf{Consequence:} at this frame budget the extractor is not
              the limiting stage of accuracy
      \end{itemize}
    \column{0.48\textwidth}
      \includegraphics[width=\textwidth,height=0.78\textheight,keepaspectratio]{\FIG/realistic_trajectories.png}
  \end{columns}
  \src{Solid = closer active, dashed = disabled, lighter = GPU front end}
\end{frame}

% ===================================================================== Part 3
\section{Experiment 3: cost}

\begin{frame}{3. Cost is measured on the two mapping runs, three repeats each}
  \begin{columns}[T]
    \column{0.44\textwidth}
      \footnotesize
      \begin{itemize}
        \item The four short recordings cannot support an uncertainty: one run
              per cell
        \item The mapping runs are \SI{288}{} and \SI{315}{\second} over
              \SI{374}{} and \SI{385}{\metre}, every aisle entered
        \item They are also the only recordings that reach the project's
              \SI{1.5}{\metre\per\second}: median \SI{1.49}{} and
              \SI{1.48}{\metre\per\second}
        \item Eighteen runs: ORB-SLAM3 CPU and GPU closure-free, and
              VINS-Fusion, three repeats of each
      \end{itemize}
    \column{0.56\textwidth}
      \centering
      \includegraphics[width=0.80\textwidth,height=0.66\textheight,keepaspectratio]{\FIG/scene_dataset_allsensor_realistic_map.png}\\
      {\scriptsize \repo{realistic_map}: the U-layout mapping route}
  \end{columns}
  \src{\repo{exp3-rules/mapping_resources.csv}, 18 runs}
\end{frame}

\begin{frame}{3. Neither pipeline is cheaper: it is a memory-against-cores trade}
  \centering\scriptsize\setlength{\tabcolsep}{4pt}
  \begin{tabular}{llrrrrr}
    \toprule
    Recording & Pipeline & Front [\si{\milli\second}] & Back [\si{\milli\second}]
      & Kept & CPU [\si{\percent}] & RSS [MiB] \\
    \midrule
    \repo{realistic_map} & ORB-SLAM3 CPU & 42.0 & 211.3 & 93 & 379 & 1732 \\
      & ORB-SLAM3 GPU & 36.7 & 284.2 & 96 & 332 & 1794 \\
      & VINS-Fusion & 38.1 & \good{44.3} & 91 & \bad{819} & \good{479} \\
    \addlinespace
    \repo{realistic_01_map} & ORB-SLAM3 CPU & 43.9 & 220.4 & 91 & 381 & 1816 \\
      & ORB-SLAM3 GPU & 51.9 & 405.7 & 75 & 361 & 1719 \\
      & VINS-Fusion & 38.6 & \good{44.5} & 90 & \bad{814} & \good{497} \\
    \bottomrule
  \end{tabular}
  \par\vspace{0.6em}
  \footnotesize Mean over three repeats. VINS-Fusion optimises five to nine
  times faster than ORB-SLAM3's local mapping and holds its state in a quarter
  of the memory, for more than twice the processor time: about eight cores of
  one against three and a half. Front ends are within \SI{6}{\milli\second}.
  \src{Mean of three repeats per configuration}
\end{frame}

\begin{frame}{3. The I-layout front-end comparison is confounded with heat}
  \footnotesize
  \begin{itemize}
    \item On the U recording the GPU front end behaves as the single runs
          showed: \SI{37}{} against \SI{42}{\milli\second},
          \SI{96}{} against \SI{93}{\percent} kept
    \item On the I recording the sign reverses: \SI{52}{} against
          \SI{44}{\milli\second}, \SI{75}{} against \SI{91}{\percent} kept
    \item \textbf{But} each configuration was run as its own block: six GPU
          runs 21:04--22:04, six CPU runs 23:27--00:18, six VINS runs the next
          midday, on a laptop that throttles at \SI{90}{\celsius}
    \item VINS-Fusion reads the \emph{same} on both recordings
          (\SI{38.1}{} and \SI{38.6}{\milli\second} front end), so the
          recordings themselves do not impose a difference of that size
    \item \bad{Quote the U figures} until the three configurations are
          repeated interleaved from a cool start
  \end{itemize}
  \src{Run start times from each run's \repo{run_metadata.csv}}
\end{frame}

% ===================================================================== Part 4
\section{Experiment 1: real baseline}

\begin{frame}{4. In simulation the EKF loses to VINS once both share a noise model}
  \begin{columns}[T]
    \column{0.46\textwidth}
      \footnotesize
      \begin{itemize}
        \item Design decision: the EKF baseline now runs on the inertial noise
              the \emph{visual} configurations declare, not its own measured
              model
        \item Reason: the two families must share one hand-set model before
              their accuracies compare at all
        \item At that model the EKF holds
              \SIrange{0.68}{2.80}{\metre} and loses to VINS-Fusion
              (\SIrange{0.22}{0.66}{\metre}) on all four
        \item The same filter on its measured model holds
              \SIrange{0.21}{0.58}{\metre} and would beat VINS on three of four
        \item \textbf{So the whole gap is the noise model, none of it the
              architecture}
      \end{itemize}
    \column{0.54\textwidth}
      \includegraphics[width=\textwidth,height=0.80\textheight,keepaspectratio]{\FIG/proprio_trajectories_grid_realistic.png}
  \end{columns}
  \src{Configured gyroscope density is 5--7$\times$ the measured value}
\end{frame}

\begin{frame}{4. The real robot now has a non-visual baseline, and vision does not win}
  \begin{columns}[T]
    \column{0.44\textwidth}
      \footnotesize
      \begin{itemize}
        \item The rig publishes no metric wheel speed and no joint states; its
              steering feedback is a two-valued command echo, so no bicycle
              model and no EKF
        \item What works: speed feedback $\times$ one scale $k$, heading from
              the gyroscope with its rest bias removed
        \item $k$, the gyro \emph{axis} and its sign are fitted on
              \repo{realv7_000} only, then held
        \item The yaw rate is on the sensor's \textbf{y} axis; its z axis
              integrates to zero over a \SI{6.3}{\radian} heading sweep
      \end{itemize}
    \column{0.56\textwidth}
      \includegraphics[width=\textwidth,height=0.52\textheight,keepaspectratio]{\OUTR/exp1-wheel/wheel_paths_light.png}
      \par\centering\scriptsize\setlength{\tabcolsep}{3pt}
      \begin{tabular}{llrr}
        \toprule
        Recording & Role & Wheel+gyro & Best ORB-SLAM3 \\
        \midrule
        \repo{realv7_000} & fit & 0.06 & 0.25 \\
        \repo{realv7_001} & held & \good{0.26} & 0.44 \\
        \repo{realv7_002} & held & \good{0.22} & 0.47 \\
        \repo{realv6_000} & held & 0.46 & 0.46 \\
        \bottomrule
      \end{tabular}
  \end{columns}
  \src{Planar ATE in metres; lever arm zero on both sides}
\end{frame}

\begin{frame}{4. What bounds that comparison}
  \footnotesize
  \begin{itemize}
    \item \bad{One tuned constant.} $k$ is fitted against the same tracker the
          baseline is scored on, albeit on a different recording. The fit
          recording's \SI{0.06}{\metre} is a calibration, not a result
    \item The fitted $k=0.325$ agrees with the plain path-length ratio of each
          \repo{realv7} bag (0.323--0.329), which is a weak independent check
    \item Path ratios come out within \SI{4}{\percent} of one, where the visual
          runs read \SIrange{5}{15}{\percent} long
    \item Routes are \SIrange{22}{37}{\metre} at \SI{0.2}{\metre\per\second};
          gyroscope drift grows with route length and vision's need not
    \item The marker-to-IMU lever arm is unmeasured on both sides
    \item \textbf{Fix:} a measured scale from wheel circumference and gear
          ratio, and the steering angle recorded as a measurement, which would
          also restore the EKF with the \SI{0.60}{\metre} wheelbase
  \end{itemize}
  \src{\repo{script/real_wheel_baseline.py}, \repo{output/compare/real/exp1-wheel/}}
\end{frame}

% ===================================================================== Part 5
\section{Experiment 2: detector}

\begin{frame}{5. The four-class checkpoint makes the pallet jacks scorable}
  \begin{columns}[T]
    \column{0.52\textwidth}
      \footnotesize
      \begin{itemize}
        \item \repo{weight/runs_1/.../best.pt}: bin, box, bucket,
              \textbf{forklift}; the fourth fires on the pallet jacks, which
              are the moving objects
        \item Scored on all four realistic recordings, \num{12698} frames
        \item Pooled precision \textbf{0.747}; class F1 \textbf{0.994} given a
              match
        \item Recall \textbf{0.502} over the boxes the exemptions keep,
              \textbf{0.288} over everything in view
        \item Moving-object recall \textbf{0.288} --- the ceiling on any
              filtering benefit downstream
        \item \bad{The walking people are simulator actors}: no pose, no mesh,
              outside the class set, so unscorable \emph{and} unmaskable
      \end{itemize}
    \column{0.48\textwidth}
      \includegraphics[width=\textwidth,height=0.72\textheight,keepaspectratio]{\FIG/yolo_fig_size.png}
  \end{columns}
  \src{\repo{output/yolo_eval/*_4class/}; confidence 0.35, matching IoU 0.5}
\end{frame}

\begin{frame}{5. The miss is a size effect; the false positives are near misses}
  \begin{columns}[T]
    \column{0.40\textwidth}
      \footnotesize
      \begin{itemize}
        \item Frame-level recall rises monotonically with truth-box side:
              0.02 at 24--32\,px to 0.64 above 192\,px
        \item The moving pallet jacks lie on the \emph{same} curve as the
              parked ones, so 0.29 is a distance effect, not a motion effect
        \item Of the unmatched predictions,
              \SIrange{74}{79}{\percent} overlap the right object below the
              IoU threshold; only \SIrange{8}{9}{\percent} land on nothing
        \item A filter that masks by box does not care which
        \item Misses come as flicker: \SIrange{39}{47}{\percent} of gaps are a
              single frame
      \end{itemize}
    \column{0.60\textwidth}
      \includegraphics[width=\textwidth,height=0.80\textheight,keepaspectratio]{\FIG/yolo_fig_fp_anatomy.png}
  \end{columns}
  \src{\repo{output/compare/simulation/exp2-insights/}}
\end{frame}

% ===================================================================== Part 6
\section{Experiment 4: maps}

\begin{frame}{6. All four realistic maps fail the criterion fixed before scoring}
  \centering\scriptsize\setlength{\tabcolsep}{5pt}
  \begin{tabular}{llrrrr}
    \toprule
    Database & Layout, build & Precision & Recall seen & False obstacles & Drivable \\
    \midrule
    \repo{map_realistic} & U, October & \bad{81.2} & \bad{82.5} & \bad{3.84} & 96.1 \\
    \repo{..._realistic_map} & U, November & \bad{77.8} & \bad{74.9} & \bad{5.00} & 94.9 \\
    \repo{map_realistic_01} & I, October & \bad{84.8} & \bad{81.4} & \bad{4.39} & 94.9 \\
    \repo{..._realistic_01_map} & I, November & \bad{82.2} & \bad{77.1} & \bad{10.19} & 89.4 \\
    \midrule
    Criterion & fixed before scoring & $\ge 90$ & $\ge 85$ & $\le 1$ & --- \\
    \bottomrule
  \end{tabular}
  \par\vspace{0.6em}
  \footnotesize All values in percent, probed with a recording that did not
  build the map. Two configurations $\times$ two layouts, one build each, so
  the parameter differences below are signatures rather than ablations.
  \src{\repo{exp4-realistic-maps/realistic_map_metrics.csv}}
\end{frame}

\begin{frame}{6. The November configuration is worse, through three parameters}
  \footnotesize
  \begin{itemize}
    \item The November builds differ from the October ones in exactly three
          RTAB-Map parameters, and each leaves its own signature:
    \begin{itemize}\footnotesize
      \item lower detection rate $\rightarrow$ fewer nodes, fewer views per cell
      \item longer grid range $\rightarrow$ phantom cells along the west wall
      \item higher ground-height cut $\rightarrow$ obstacles in the aisles, and
            on the I layout the blocking phantoms appear \emph{mid-route}
            (median first node 27, against 1 everywhere else)
    \end{itemize}
    \item The false obstacles are \textbf{not} thickened shelving: they are
          phantoms laid during the inertial initialisation transient at the
          spawn, plus wall bands displaced about \SI{0.5}{\metre}
    \item \textbf{The correction the data predict:} discard the cells of the
          first few nodes, or start mapping after the transient settles
    \item The three parameters changed together, so which one dominates is
          untested
  \end{itemize}
  \src{\repo{script/exp4_insights.py --realistic}, Insight 6}
\end{frame}

% ===================================================================== Part 7
\section{Report and next steps}

\begin{frame}{7. The report is rebuilt on the fixed dataset set}
  \footnotesize
  \begin{itemize}
    \item \textbf{Chapter 3} purged of the retired campaign: nine fragments of
          the old report dropped and rewritten from current measurements ---
          registry, worlds, sensor fidelity, the four experiments' methods,
          status and deviations, real rig
    \item Every number in it re-measured on the recordings in use: wheel slip
          under \SI{0.4}{\percent}, gyroscope noise 5--7$\times$ below the
          configured value, the inertial orientation field confirmed to be
          ground truth
    \item \textbf{Chapter 4} slots 1--3 are now generated from the insight
          notes by \repo{script/build_chapter4.py}; edit the note, regenerate
    \item \textbf{Chapter 5} rewritten; the two hypotheses the project is named
          for still have no real-robot evidence
    \item 129 pages, no undefined references, four overfull lines, all
          pre-existing
  \end{itemize}
  \src{\repo{script/build_final_report.sh}}
\end{frame}

\begin{frame}{7. What is open, in the order it should be closed}
  \footnotesize
  \begin{enumerate}
    \item \textbf{Give Experiment~3 its test.} Both pipelines filtered and
          unfiltered on \repo{realistic_002} and \repo{_003}, paired under the
          admissibility gates, with the closer \emph{disabled} on the I layout
          so the false closure does not dominate both arms
    \item \textbf{Bound the loop closer.} A maximum correction against the
          inertial drift bound, tested on \repo{realistic_001} where the
          \SI{14}{\metre} closure is reproducible in both front ends
    \item \textbf{Repeat the cost runs interleaved} from a cool start, so the
          I-layout front-end comparison is readable
    \item \textbf{Real robot:} exercise the fork's initialisation switches,
          re-run VINS with the corrected gravity norm, measure the
          marker-to-IMU lever arm, measure the wheel-speed scale
    \item \textbf{Detector:} a person class with projectable geometry, or an
          annotated frame subset, so the agents the filter exists for are
          scorable
    \item \textbf{Maps:} start mapping after the inertial transient, and
          reconsider the November range and ground-height parameters
  \end{enumerate}
\end{frame}

\begin{frame}{Summary}
  \footnotesize
  \begin{itemize}
    \item The report is now built from eleven recordings and nothing else;
          everything retired is logged, not deleted
    \item \textbf{ORB-SLAM3's error in simulation is its loop closer}, in both
          directions: a false \SI{14}{\metre} closure on the I layout's
          repeated aisles, and a vertical drift on the same layout that only a
          closure removes. The front end changes retention, not accuracy
    \item \textbf{VINS-Fusion} is the steadier estimator in simulation and
          costs a quarter of the memory for twice the cores, but still has no
          valid run on the real robot
    \item \textbf{Vision has not earned its place over dead reckoning} in
          either domain, on routes of this length
    \item \textbf{The detector finds under a third of the moving objects},
          which bounds any filtering benefit before the filter is tested
    \item \textbf{All four realistic maps fail} the map-quality criterion, and
          the cause is the initialisation transient, not the shelving
  \end{itemize}
\end{frame}

\end{document}
